\documentclass[11pt]{article}

\usepackage[T1]{fontenc}
\usepackage[margin=1in]{geometry}
\usepackage{amsmath, amssymb, amsthm}
\usepackage{enumitem}
\usepackage{hyperref}
\usepackage{booktabs}
\usepackage{listings}
\usepackage{xcolor}

\lstset{
  basicstyle=\ttfamily\small,
  frame=single,
  breaklines=true,
  columns=fullflexible,
  mathescape=true
}

\theoremstyle{definition}
\newtheorem{theorem}{Theorem}[section]
\newtheorem{definition}[theorem]{Definition}
\newtheorem{example}[theorem]{Example}

\title{CS 250 --- Intermezzo: The Algebra of Types}
\author{}
\date{}

\begin{document}
\maketitle

Something strange has been happening throughout this course, and it's time we talked about it.

In Module 1 we learned that $|A \times B| = |A| \cdot |B|$---the cardinality of a Cartesian product is the product of the cardinalities. In the sets intermezzo we learned that $|B^A| = |B|^{|A|}$---function spaces have exponential cardinality. And in Module 10 we defined data types using BNF, building lists and trees out of simpler pieces.

These facts live in different chapters, but they are secretly the same idea. Types obey the laws of algebra. Not metaphorically, not ``sort of''---they obey the \emph{actual} laws of high-school algebra, the ones with $+$ and $\times$ and exponents. Once you see this, a lot of things click into place.

\section{Types are like numbers}

Let's build the dictionary.

\medskip
\noindent\textbf{Product types are multiplication.} A pair $(a, b)$ where $a \in A$ and $b \in B$ has type $A \times B$. How many inhabitants does $A \times B$ have? For each choice of $a$ (there are $|A|$ many) and each choice of $b$ (there are $|B|$ many), we get a distinct pair. So $|A \times B| = |A| \cdot |B|$. A tuple is multiplication.

\medskip
\noindent\textbf{Sum types are addition.} Suppose we have a type $A + B$ that means ``either a value of type $A$, tagged with \texttt{Left}, or a value of type $B$, tagged with \texttt{Right}.'' In Haskell this is \texttt{Either A B}; in Python you might use a tagged union or a class hierarchy. How many inhabitants? Every element of $A$ gives one (tagged \texttt{Left}), and every element of $B$ gives another (tagged \texttt{Right}). So $|A + B| = |A| + |B|$. A tagged union is addition.

\medskip
\noindent\textbf{Function types are exponentiation.} A function $f : A \to B$ assigns to each element of $A$ an element of $B$. That's $|B|$ choices for each of the $|A|$ inputs, so $|A \to B| = |B|^{|A|}$. We saw this in the sets intermezzo---it's why the function space is written $B^A$. A function is exponentiation.

\medskip
\noindent\textbf{The unit type is 1.} Let $1$ denote a type with exactly one element (in Haskell, \texttt{()};  in Python, \texttt{None}). Then $A \times 1 \cong A$---a pair of an $A$ and a ``nothing'' is just an $A$, since the second component carries no information. And $1^A \cong 1$---there's only one function from $A$ to a one-element set (the function that sends everything to the single element).

\medskip
\noindent\textbf{The empty type is 0.} Let $0$ denote a type with \emph{no} elements (in Haskell, \texttt{Void}; in Python, a class you can never instantiate). Then $A \times 0 \cong 0$---you can't build a pair if the second component has no possible values. And $A + 0 \cong A$---``either an $A$ or an impossible thing'' is just an $A$. And for any nonempty $A$, $0^A \cong 0$---there are no functions from $A$ to an empty set (you'd have to send the elements of $A$ somewhere, but there's nowhere to send them).

\medskip
Look at these equations again. $A \times 1 = A$. $A + 0 = A$. $A \times 0 = 0$. Those are the rules of arithmetic you learned in grade school, but now they're facts about data types. Types form an algebra, and it's the same algebra you already know.

\section{The laws of algebra hold}

Let's verify that the familiar identities from arithmetic actually hold for types. In each case, ``$\cong$'' means ``there's a bijection between the two types''---they have the same number of inhabitants, and you can convert back and forth without losing information.

\medskip
\noindent\textbf{Multiplicative identity:} $A \times 1 \cong A$. A pair where one component is trivial is just the other component. In Haskell: \texttt{(a, ())} is isomorphic to \texttt{a}.

\medskip
\noindent\textbf{Additive identity:} $A + 0 \cong A$. \texttt{Either A Void} is isomorphic to \texttt{A}---the \texttt{Right} case is impossible (you can't construct a \texttt{Void}), so you're always in the \texttt{Left} case.

\medskip
\noindent\textbf{Distributivity:} $A \times (B + C) \cong (A \times B) + (A \times C)$. A pair of an $A$ with ``either a $B$ or a $C$'' is the same as ``either a pair of $A$ and $B$, or a pair of $A$ and $C$.'' This is the distributive law $a(b + c) = ab + ac$, in type form.

\medskip
\noindent\textbf{Exponent of a product:} $A \to (B \times C) \cong (A \to B) \times (A \to C)$. A function that returns a pair is the same as a pair of functions. If I ask you for a machine that takes an input and produces two outputs, you can either build one machine with two output slots or build two separate machines---same thing.

\medskip
\noindent\textbf{Exponent of a sum:} $(A + B) \to C \cong (A \to C) \times (B \to C)$. A function from ``either an $A$ or a $B$'' is the same as a pair of functions: one that handles $A$ inputs and one that handles $B$ inputs. In Python:

\begin{lstlisting}[language=Python]
# A function from Either is a pair of handlers:
def handle_either(x):
    if isinstance(x, Left):
        return handle_a(x.value)   # the A -> C part
    else:
        return handle_b(x.value)   # the B -> C part
\end{lstlisting}

Wait a moment. A function from a sum is a pair of handlers---one for each case. Isn't that exactly \emph{proof by cases}? In Module 3 we learned the disjunction elimination rule: if you want to prove something from $A \vee B$, you handle the $A$ case and handle the $B$ case separately. That's the logical version of $(A + B) \to C \cong (A \to C) \times (B \to C)$---or-elimination \emph{is} the fact that a function from a sum is a pair of functions. The algebra is telling us something about logic.

\section{Data types as polynomials}

Here is where things get a little bit magical.

In Module 10 we defined lists with a BNF:

\begin{lstlisting}
List A := Nil | Cons(A, List A)
\end{lstlisting}

In our algebraic notation, \texttt{Nil} is the unit type $1$ (one way to be empty) and \texttt{Cons(A, List A)} is a product $A \times L$ where $L$ is the type of lists. The \texttt{|} in the BNF is a sum. So the definition of lists is:
\[
L = 1 + A \cdot L
\]

That's a polynomial equation! Let's be reckless and solve it like one. Rearranging:
\[
L - A \cdot L = 1 \qquad \Longrightarrow \qquad L(1 - A) = 1 \qquad \Longrightarrow \qquad L = \frac{1}{1 - A}
\]

Now, you might be having a reasonable objection right now, which is: what does it even \emph{mean} to divide types or subtract types? Fair enough. But let's push ahead and see what happens. From Module 9, we know the formula for a geometric series:
\[
\frac{1}{1 - x} = 1 + x + x^2 + x^3 + \cdots
\]

So if we just blindly expand:
\[
L = 1 + A + A^2 + A^3 + A^4 + \cdots
\]

And what is this saying? A list of $A$'s is:
\begin{itemize}
  \item $1$: the empty list (one way), or
  \item $A$: a single element (one $A$), or
  \item $A^2 = A \times A$: two elements (a pair of $A$'s), or
  \item $A^3 = A \times A \times A$: three elements, or
  \item $\ldots$and so on forever.
\end{itemize}

That's \emph{exactly right}. A list is either empty, or has one element, or has two, or three, or any finite number of elements. The algebra of types---this insane-looking manipulation where we divided by $(1 - A)$ as if types were numbers---told us the truth. The geometric series from Module 9 is secretly a statement about data structures.

\medskip
Let's try another one. Binary trees from Module 10 were defined as:

\begin{lstlisting}
Tree := Leaf | Node(Tree, Tree)
\end{lstlisting}

In algebraic notation:
\[
T = 1 + T^2
\]

This is a quadratic equation! Rearranging: $T^2 - T + 1 = 0$, and the quadratic formula gives
\[
T = \frac{1 \pm \sqrt{1 - 4}}{2} = \frac{1 \pm \sqrt{-3}}{2}
\]

which involves complex numbers and seems completely unhinged. But remarkably, if you take the right root and expand it as a power series, you get the \emph{Catalan numbers}---which are exactly the number of distinct binary trees with $n$ internal nodes. The algebra of types knows about combinatorics.

We won't push this further here, but the point is: these aren't parlor tricks. There is a rigorous mathematical theory---the theory of \emph{combinatorial species}---that makes all of this precise. The equation $L = 1 + A \cdot L$ is a legitimate algebraic equation in that theory, the ``division'' is a real operation, and the geometric series expansion is a theorem, not a coincidence.

What we're seeing is a remarkable convergence: the sets intermezzo (cardinality of function spaces), Module 9 (geometric series), and Module 10 (BNF definitions of data types) are all aspects of a single algebraic framework. Types are numbers, data definitions are polynomial equations, and the power series you learned in Module 9 enumerate the shapes of data structures.

\section{The Curry--Howard connection}

There's one more layer. In the Curry--Howard intermezzo we learned that types correspond to propositions: a type is ``true'' if it's inhabited (has at least one value) and ``false'' if it's empty. Under this correspondence:

\begin{center}
\begin{tabular}{lll}
\toprule
\textbf{Algebra} & \textbf{Types} & \textbf{Logic} \\
\midrule
$A \cdot B$ & $A \times B$ (product) & $A \wedge B$ (conjunction) \\
$A + B$ & $A + B$ (sum) & $A \vee B$ (disjunction) \\
$B^A$ & $A \to B$ (function) & $A \to B$ (implication) \\
$1$ & Unit & $\top$ (true) \\
$0$ & Void & $\bot$ (false) \\
\bottomrule
\end{tabular}
\end{center}

So the algebraic identities we verified in Section 2 are also \emph{logical tautologies}:
\begin{itemize}
  \item $A \times 1 \cong A$ is $A \wedge \top \equiv A$ (``$A$ and true'' is just $A$)
  \item $A + 0 \cong A$ is $A \vee \bot \equiv A$ (``$A$ or false'' is just $A$)
  \item $A \times (B + C) \cong (A \times B) + (A \times C)$ is the distribution of $\wedge$ over $\vee$
  \item $(A + B) \to C \cong (A \to C) \times (B \to C)$ is or-elimination
  \item $A \times 0 \cong 0$ is $A \wedge \bot \equiv \bot$ (``$A$ and false'' is false)
\end{itemize}

Three parallel worlds---arithmetic, type theory, and logic---and they're all obeying the same laws. The algebra of numbers, the algebra of types, and the algebra of propositions are three views of one structure.

That structure has a name, by the way. It's called a \emph{rig}\index{rig} (a ``ring without negation''---because you can't subtract types or negate propositions in general). But you don't need to know the name to appreciate the fact: the math you already know, from all across this course, fits together like a puzzle. Cardinality is counting. Types are numbers. BNF definitions are polynomial equations. Proofs are programs. And the geometric series from Module 9 secretly counts the shapes of lists.

\section{Exercises}

\bigskip
\textbf{Exercise 1.}\label{ex:ta:1} Using the algebra of types, simplify the type $(A \to B) \times (A \to C)$. What single type is it isomorphic to? What does this isomorphism correspond to as a logical proposition (under the Curry--Howard correspondence)?

\bigskip
\textbf{Exercise 2.}\label{ex:ta:2} The BNF for natural numbers is \texttt{Nat = Zero | Succ(Nat)}, which in our algebraic notation is $N = 1 + N$. Write this as a ``polynomial equation'' and ``solve'' for $N$ by expanding as a series. Does the series $1 + 1 + 1 + 1 + \cdots$ make sense as a description of the natural numbers? (Hint: what does $1^n = 1$ mean for the $n$th term?)

\bigskip
\textbf{Exercise 3.}\label{ex:ta:3} (Distributivity, programmatically.) Use the algebra of types to simplify $(A + B) \times (C + D)$. What single ``polynomial'' do you get? Convert your answer back into a Python data structure: write a class hierarchy or tagged union with one constructor for each summand of the result. What logical tautology does this correspond to under Curry--Howard?

\bigskip
\textbf{Exercise 4.}\label{ex:ta:4} (Solving for non-empty lists.) Define non-empty lists by the BNF
\[
  \texttt{NEList A} = A \;\mid\; A \times \texttt{NEList A}
\]
that is, ``either a single $A$, or an $A$ followed by a non-empty list of $A$.'' Translate this to a polynomial equation in $L$ (where $L = \texttt{NEList A}$), solve for $L$ algebraically, and expand the result as a series. What kinds of values appear in the series, and how does that match the structure of a non-empty list?

\bigskip
\textbf{Exercise 5.}\label{ex:ta:5} (The \texttt{Maybe} type.) The type \texttt{Maybe A} (in Haskell) is defined as \texttt{Just A | Nothing}. In algebraic notation: $\texttt{Maybe } A = A + 1$.
\begin{enumerate}[label=(\alph*)]
  \item What is $|\texttt{Maybe } A|$ in terms of $|A|$?
  \item What is $|\texttt{Maybe (Maybe } A\texttt{)}|$ in terms of $|A|$?
  \item Express \texttt{Maybe (Maybe A)} as an algebraic expression and simplify. Why does the answer match the cardinality formula in (b)?
\end{enumerate}

\bigskip
\textbf{Exercise 6.}\label{ex:ta:6} (A function that returns a pair is a pair of functions.) Verify the isomorphism
\[
  A \to (B \times C) \;\cong\; (A \to B) \times (A \to C)
\]
by writing two Python functions: one that converts the left side into the right side, and one that converts the right side into the left side. Verify (informally, by trying out a few examples) that they are inverses of each other. What does this say in pure-arithmetic terms about the equation $(b \cdot c)^a = b^a \cdot c^a$?

\section*{Further reading}

If the algebra of types appeals to you, here is where to go next. The connection to combinatorics is deep and beautiful.

\begin{itemize}
  \item Conor McBride, ``The Derivative of a Regular Type is its Type of One-Hole Contexts,'' unpublished manuscript (2001), available on the author's homepage at \texttt{strictlypositive.org}. A short, witty, surprisingly elementary derivation of why \emph{differentiating} an algebraic data type literally gives you its zipper. Connects high-school calculus to programming in a way undergraduates love.
  \item Brent Yorgey, ``Species and Functors and Types, Oh My!'' \emph{Haskell Symposium 2010}. Written explicitly as an accessible bridge between Joyal's combinatorial species and the algebraic data types programmers already know. The most undergraduate-friendly entry point to species I know of.
  \item Philippe Flajolet and Robert Sedgewick, \emph{\href{https://ac.cs.princeton.edu/home/}{Analytic Combinatorics}}, Cambridge University Press (2009), Chapter I (``Symbolic Methods''). The full book PDF is freely available on Sedgewick's Princeton homepage. Chapter I introduces the symbolic method---sums, products, sequences as type combinators with generating functions---at exactly the level a strong undergrad can handle.
\end{itemize}

\end{document}
